% Owner: TV3 (P5.3b) -- Reviewer: TV1
\section{Conclusion}
\label{sec:conclusion}

We compared Minimax, alpha-beta pruning, and MCTS-UCT on $3\times 3$ Tic-Tac-Toe and on Gomoku, and extended the study to Xiangqi. On $3\times 3$, alpha-beta returns the Minimax value while visiting $96.7$--$98.7\%$ fewer nodes, and move ordering alone changes its cost by a factor of $2.5$ (RQ1). MCTS with $c = \sqrt{2}/2$ stops losing to every opponent strategy only at $I^{*} = 2\,000$ iterations per move, although it already draws every game against perfect Minimax at $1\,000$ iterations (RQ3). The exploration constant matters only when the budget is small: on $9\times 9$ Gomoku with $100$ iterations, the win rate falls from $96.5\%$ to $88.5\%$ as $c$ grows from $0.5$ to $2$ (RQ4). Choosing the final move by $N$ or by $W/N$ makes little difference in play (RQ5), while the form of the exploration term does: bonuses that decay as $1/N_i$ are clearly worse, and a heuristic prior is clearly better. On Gomoku, alpha-beta depends above all on move ordering, iterative deepening, and the candidate restriction (RQ2), and it beats MCTS in $594$ of $600$ games at equal time (RQ6). Informed rollouts help MCTS at equal iterations but only the truncated rollout L4 still helps at equal time, which shows that knowledge must be cheap compared with the simulations it replaces. Natural next steps are to learn the evaluation and the prior from self-play, as AlphaGo Zero does, and to implement the search in a compiled language.
